\section{Introduction}
In recent years, the demand for intelligent vehicle monitoring systems has grown significantly, driven by the need to enhance driver and passenger safety, improve fuel efficiency, reduce emissions, and enable real-time diagnostics. With the global emphasis on climate change and the increasing integration of automation in the automotive sector, monitoring systems have evolved from simple sensor loggers to intelligent platforms capable of advanced data processing and analysis.

One of the key aspects of such monitoring is the ability to retrieve and interpret diagnostic and performance-related data from vehicles. This is crucial not only for preventive maintenance but also for post-accident forensic analysis and insurance claims. Recognizing these needs, we developed the \textbf{Automotive Diagnostic and Activity Monitoring (ADAM)} system, which utilizes the vehicle's On-Board Diagnostics (OBD-II) interface in conjunction with multiple sensing modules to create a holistic monitoring ecosystem.

\textbf{The OBD-II standard} was introduced in 1996 by the Society of Automotive Engineers (SAE) and mandated by the U.S. Environmental Protection Agency (EPA) to regulate vehicle emissions. It consists of a 16-pin standardized connector found in all modern light-duty vehicles, serving as a critical access point for monitoring emission-related data and overall vehicle health. This interface provides standardized access to sensor data, engine diagnostics, and fault codes from the vehicle’s Electronic Control Unit (ECU)—the central controller that governs all operational aspects of a vehicle.

\subsection{OBD-II and Parameter Identifiers (PIDs)}
The OBD-II standard supports several communication protocols, including:
\begin{itemize}
    \item ISO 15765-4 (CAN-BUS) – widely adopted since 2008
    \item ISO 14230-4 (KWP2000)
    \item ISO 9141-2
    \item SAE J1850 VPW
    \item SAE J1850 PWM
\end{itemize}

This study focuses on the ISO 15765-4 (Controller Area Network – CAN) protocol, which is now the default protocol in most vehicles due to its robustness and speed. OBD-II communication occurs via hexadecimal-coded requests sent to the ECU, which responds with encoded data or Diagnostic Trouble Codes (DTCs). These codes are categorized into:

\begin{itemize}
    \item \textbf{DTCs (Diagnostic Trouble Codes)} – Identify faults and malfunctions in vehicle subsystems.
    \item \textbf{PIDs (Parameter Identifiers)} – Used to request live sensor data and real-time vehicle parameters.
\end{itemize}

The following table lists commonly used PIDs in OBD-II diagnostics, along with their descriptions and typical byte responses:

\begin{table}[h]
\centering
\caption{Common OBD-II Parameter Identifiers (PIDs)}
\begin{tabular}{|c|l|l|}
\hline
\textbf{PID} & \textbf{Parameter} & \textbf{Description / Units} \\
\hline
0C & Engine RPM & Revolutions per minute (2 bytes) \\
0D & Vehicle Speed & Speed in km/h (1 byte) \\
05 & Engine Coolant Temp & Temperature in °C (1 byte) \\
0F & Intake Air Temp & Temperature in °C (1 byte) \\
11 & Throttle Position & Percentage (1 byte) \\
10 & MAF Air Flow Rate & Grams/sec (2 bytes) \\
2F & Fuel Level Input & Percentage (1 byte) \\
42 & Control Module Voltage & Voltage in volts (2 bytes) \\
\hline
\end{tabular}
\label{tab:obd_pids}
\end{table}

By utilizing these PIDs, ADAM is capable of extracting a rich set of telemetry from the vehicle, which is further processed using embedded logic and machine learning algorithms for deeper insights into engine health, fuel efficiency, driver behavior, and predictive maintenance.



% \section{Introduction}
% Over the last couple of years, vehicle monitoring has been a necessary measure to make vehicles and drivers safe, cut down on emissions to tackle climate change , and provide optimal performance. The significance of protecting human life and the necessity to obtain genuine information at the time of accidents for forensic analysis reinforce the importance of efficient monitoring systems. To serve this purpose, we created a vehicle monitoring logger on the basis of On-Board Diagnostics (OBD-II) port. The OBD-II port is a 16-pin standardized connector found in every contemporary vehicle,developed in the United States Of America in 1996, by the  Society of Automotive Engineers (SAE). It was imposed to enable regulation  of vehicle emission so as to ensure Environmental Protection Agency(EPA) standards are met. It is a diagnostic port used to monitor real-time readings of vehicle sensors and the Diagnostic Trouble Codes (DTCs) from the Electronic Control Unit (ECU), making it a valuable resource for monitoring vehicle performance.Ecu is a the main controller of the vehicle where all controls and actions are done.

% OBD II port supports communication protocols such as ISO15765-4 (CAN-BUS), ISO14230-4 (KWP2000), ISO 9141-2, SAE J1850 VPW, and SAE J1850 PWM.  This study is centered on the ISO 15765-4 (CAN-BUS) protocol that has been required on vehicles manufactured from 1996.
% The message request is send in a hexadecimal to the protocols and it gets interpreted according  to one of the five OBD II protocols and sends back a hexadecimal response.Obd II uses two types of codes to request ECU data that are Diagnostic Trouble Code (DTCs) and Parameter Identifiers (PIDs).DTC are used to diagnostic malfunctions in various subsystem of the vehicle and PIDs are used to measure real time parameters 
